% Lösungen Einheit 2 von 2 – Einheit 2 (zusammenbau v0.1)
\einheitenkopf{Einheit 2 von 2 · Einheit 2}

\erg{9}{a) ja – die Höhen ergeben zusammen eins \quad b) $P(X \le 2) = 0{,}3 + 0{,}4 = 0{,}7$ \quad c) $Y$ gehört zu III, denn nur III ist schief mit dem Schwerpunkt bei vielen Treffern. $X$ gehört zu II: Die Summe $2$ entsteht aus drei Paaren, die Summe $0$ aus einem, also $P(X = 2) : P(X = 0) = 3$ – bei I ist das Verhältnis nur $1{,}5$.}
\erg{10}{a) Fehler: den Einzelwert als kumulierte Wahrscheinlichkeit angegeben. Richtig: $P(X \le 2) = 0{,}4 + 0{,}2 = 0{,}6$ \quad b) Symmetrisch heißt: Werte links und rechts der Mitte haben paarweise dieselbe Wahrscheinlichkeit, also tragen beide Seiten gleich viel vom Rest. \quad c) $P(X \ge 500) = 0{,}25 + 0{,}5 = 0{,}75$ \quad d) Säulen der Höhen $\frac{1}{8}$, $\frac{3}{8}$, $\frac{3}{8}$, $\frac{1}{8}$ über $0$ bis $3$.}

10d)

\saeulen{0/0.125,1/0.375,2/0.375,3/0.125}{0.5}{0.125}{$P(X = k)$}

